% !TeX root = master.tex
% !TeX spellcheck = de_DE

\appendix
\chapter{Anhang}
\label{sec:anhang}
\raggedbottom

\begin{minted}[breaklines]{tcl}
    set_property -dict {PACKAGE_PIN W5  IOSTANDARD LVCMOS33 } [get_ports i_clk];
    set_property -dict {PACKAGE_PIN U18  IOSTANDARD LVCMOS33 } [get_ports i_rst_n];
    set_property -dict {PACKAGE_PIN U16  IOSTANDARD LVCMOS33 } [get_ports q];
    #set_property -dict {PACKAGE_PIN H5  IOSTANDARD LVCMOS33 } [get_ports q]

    set_property CFGBVS VCCO [current_design];
    set_property CONFIG_VOLTAGE 3.3 [current_design];

    create_clock -add -name sys_clk_pin -period 10.00 -waveform {0 5} [get_ports i_clk];
\end{minted}
\captionof{listing}[Constraint-Datei für Artix-7 FPGA]{Constraint-Datei mit Mapping für LED, Reset-Knopf und Clock.}
\label{lst:constraint-datei}

\newpage

\begin{minted}{python}
    #!/usr/bin/env python3

    import subprocess
    import time
    import csv
    import statistics
    import os
    import signal

    cmd = [
        "stdbuf",
        "-oL",
        "fusesoc",
        "run",
        "--target=verilator_tb",
        "award-winning:serv:servant",
        "--firmware=collatzBareMetal.hex",
        "--memsize=1048576",
    ]

    anzahl_messungen = 100
    laufzeiten = []
    csv_datei = "laufzeitenBareMetalAwardWinningServ.csv"

    process = subprocess.Popen(
        cmd,
        stdout=subprocess.PIPE,
        stderr=subprocess.STDOUT,
        text=True,
        start_new_session=True,
    )
    print(f"\n=== Starte Zeitmessungen ===")
    start_time = None
    try:
        for line in process.stdout:
            line = line.rstrip("\r\n")
            # print(line)
            if "output q is" not in line:
                continue
            timestamp = time.perf_counter()
            if start_time is None:
                start_time = timestamp
                continue
            laufzeit = timestamp - start_time
            start_time = timestamp
            laufzeiten.append(laufzeit)
            # print(f"Laufzeit: {laufzeit:.6f} s")
            if len(laufzeiten) >= anzahl_messungen:
                print("Alle Messungen gesammelt")
                break

    # komplette Prozessgruppe beenden
    finally:
        print("Beende Simulation...")
        os.killpg(
            os.getpgid(process.pid),
            signal.SIGTERM
        )
        # falls der Prozess nicht sauber beendet wurde
        try:
            process.wait(timeout=5)
        except subprocess.TimeoutExpired:
            os.killpg(
                os.getpgid(process.pid),
                signal.SIGKILL
            )

    # CSV schreiben
    with open(csv_datei, "w", newline="") as f:
        writer = csv.writer(f)
        writer.writerow(["MessungNr", "Laufzeit_s"])
        for i, t in enumerate(laufzeiten, start=1):
            writer.writerow([i, t])

    # Statistik
    print("\n=== Statistik ===")
    print(f"Anzahl Messungen: {len(laufzeiten)}")
    print(f"Mittelwert:       {statistics.mean(laufzeiten):.6f} s")
    print(f"Minimum:          {min(laufzeiten):.6f} s")
    print(f"Maximum:          {max(laufzeiten):.6f} s")
    print(f"\nCSV gespeichert: {csv_datei}")
\end{minted}
\captionof{listing}[Python-Skript für Simulationsbenchmark]{Python-Skript, welches genutzt wurde, um Zeitmessungen in der Simulation vorzunehmen. Die resultierende .csv-Datei ist im Github-Repository \parencite{GitHub-Bachelorarbeit} zu finden.}
\label{lst:python-skript}